% GENERATED FILE. Edit canonical lessons, then run npm run generate:books.
% canonical-source-hash: fnv1a64:241043a5aeda62c9
% canonical-lessons: SA-C263-gamisyami, SA-C263-first-person-future, SA-C263-future-person-contrast, SA-C263-tomorrow-going-recall

\chapter{I will go tomorrow}
\label{ch:sa-first-person-future-going}
\hlchaptermodality{263}

\begin{chapteropening}
\textbf{By the end of this chapter:} \emph{I can distinguish will go from I will go and say that I will go tomorrow.}

Complete chapter 263 by producing a first-person future plan with a known time word.
\end{chapteropening}

\section[\texorpdfstring{\sk{गमिष्यामि} (gamiṣyāmi)}{गमिष्यामि (gamiṣyāmi)}]{\sk{गमिष्यामि} (gamiṣyāmi) — I will go}

\label{lesson:SA-C263-gamisyami}



\begin{quote}
Close the last chapter with one quick choice: several people including the speaker are the object, so say \textbf{\sk{अस्मान्}} (\emph{asmān}). Now say \textbf{\sk{गमिष्यति}} (\emph{gamiṣyati}), “will go.” That known form is the bridge to one new speaker form.
\end{quote}

\subsection*{You'll want to know: \sk{गमिष्यामि}}
\textbf{\sk{गमिष्यामि}} — \emph{gamiṣyāmi} — means “I will go.” Hear the long \textbf{-āmi} at the end. For now, learn the complete form as one usable promise.

\begin{cousinweb}
\textbf{\sk{गमिष्यामि}} comes from the root \textbf{\sk{गम्}} (\emph{gam}), “go,” with the ordinary future formation \textbf{-\sk{इष्य}-} (\emph{-iṣya-}), and ends in first-person \textbf{-\sk{आमि}} (\emph{-āmi}). Whitney describes this Sanskrit s-future and its \textbf{-iṣya- / -sya-} formation in \emph{Sanskrit Grammar} §§932–935. This lesson makes no claim beyond that transparent grammatical history.
\end{cousinweb}

\subsection*{Your turn}
Say these aloud:

\begin{itemize}
\raggedright
  \item “will go” — \emph{gamiṣyati}
  \item “I will go” — \emph{gamiṣyāmi}
  \item hear \emph{gamiṣyāmi} and answer “I will go”
\end{itemize}

\begin{tcolorbox}[breakable,colback=teal!4,colframe=teal!35!black,title={Before you move on}]
Without looking back, say “I will go.” (\textbf{\sk{गमिष्यामि}}.)
\end{tcolorbox}
\section[(first-person singular future)]{The speaker inside the future}

\label{lesson:SA-C263-first-person-future}



\begin{quote}
Contrast \textbf{\sk{गमिष्यति}}, “will go,” with \textbf{\sk{गमिष्यामि}}, “I will go.”
\end{quote}

\begin{grammarlens}[title={Grammar Lens: first-person singular}]
In the complete form \textbf{\sk{गमिष्यामि}}, the ending \textbf{-\sk{आमि}} (\emph{-āmi}) marks the speaker: “I.” The future portion comes before it. This is a \textbf{first-person singular} ramp. Learn this whole form before extending the pattern.
\end{grammarlens}

\subsection*{Your turn}
Say these aloud:

\begin{itemize}
\raggedright
  \item another person “will go” — \emph{gamiṣyati}
  \item the speaker says “I will go” — \emph{gamiṣyāmi}
  \item name the ending that carries “I” — \emph{-āmi}
\end{itemize}

\begin{tcolorbox}[breakable,colback=teal!4,colframe=teal!35!black,title={Before you move on}]
Which complete taught form means “I will go,” and which ending marks the speaker? (\textbf{\sk{गमिष्यामि}}; \textbf{-\sk{आमि}}.)
\end{tcolorbox}
\section[Practice]{Who will go?}

\label{lesson:SA-C263-future-person-contrast}



\begin{quote}
Listen for the ending before you translate.
\end{quote}

\subsection*{Your turn}
Say these aloud:

\begin{itemize}
\raggedright
  \item “he will go” — \emph{saḥ gamiṣyati}
  \item “I will go” — \emph{ahaṃ gamiṣyāmi}
  \item teacher points away — \emph{gamiṣyati}
  \item teacher points to you as speaker — \emph{gamiṣyāmi}
\end{itemize}

\begin{tcolorbox}[breakable,colback=teal!4,colframe=teal!35!black,title={Before you move on}]
Say the pair in order: “will go; I will go.” (\emph{Gamiṣyati; gamiṣyāmi.})
\end{tcolorbox}
\section[Practice]{Tomorrow, I will go}

\label{lesson:SA-C263-tomorrow-going-recall}



\begin{quote}
Recall the one word for “tomorrow.” (\textbf{\sk{श्वः}}, \emph{śvaḥ}.)
\end{quote}

\subsection*{Your turn}
Put the known time word first, then choose the speaker form.

Say these aloud:

\begin{itemize}
\raggedright
  \item “Tomorrow, I will go” — \emph{śvaḥ gamiṣyāmi}
  \item “Tomorrow, he will go” — \emph{śvaḥ saḥ gamiṣyati}
  \item change “he will go” to “I will go” — \emph{gamiṣyāmi}
\end{itemize}

\begin{tcolorbox}[breakable,colback=teal!4,colframe=teal!35!black,title={Before you move on}]
You have made a decision for tomorrow. Say only: “Tomorrow, I will go.” (\textbf{\sk{श्वः} \sk{गमिष्यामि।}})
\end{tcolorbox}
